\documentclass[../DoAn.tex]{subfiles}
\begin{document}

\begin{center}
    \Large{\textbf{TÓM TẮT NỘI DUNG ĐỒ ÁN}}\\
\end{center}
\vspace{1cm}
Robot lặn điều khiển từ xa (ROV) được ứng dụng trong khảo sát, giám sát và thao tác dưới nước, nơi tín hiệu GPS không hoạt động, tầm nhìn hạn chế và robot chịu ảnh hưởng của dòng chảy. Các giải pháp hiện nay chủ yếu gồm điều khiển thủ công, ổn định bằng cảm biến và điều khiển dựa trên thị giác máy tính. Tuy nhiên, điều khiển thủ công phụ thuộc vào kỹ năng người vận hành; cảm biến chuyên dụng có chi phí cao; còn phương pháp thị giác dễ bị tác động bởi ánh sáng và độ đục của nước. Vì vậy, đồ án lựa chọn xây dựng hệ thống phần mềm điều khiển và quản lý ROV tích hợp nhiều nguồn dữ liệu nhằm nâng cao khả năng vận hành, giám sát và lưu trữ nhiệm vụ. Giải pháp gồm phần mềm Ground Control Station (GCS) tại trạm topside và nền tảng Web. GCS cho phép điều khiển ROV bằng tay cầm, điều chỉnh tốc độ, tay gắp và đèn; đồng thời hiển thị video, hướng, độ sâu, nhiệt độ, độ ẩm và dung lượng pin. Hệ thống hỗ trợ giữ hướng, giữ độ sâu, giữ vị trí bằng DVL và giữ vị trí bằng thuật toán thị giác; ngoài ra còn ghi video, sử dụng sonar để quan sát môi trường và lưu thông tin chuyến đi như thời gian, quãng đường và quỹ đạo. Nền tảng Web tiếp nhận dữ liệu chuyến đi để lưu trữ, tra cứu và sử dụng mô hình YOLO phát hiện vật thể. Phần ROV dưới nước gồm máy tính nhúng, cảm biến và cơ cấu chấp hành chỉ là nền tảng phần cứng kết nối, không thuộc phạm vi thiết kế chính. Đóng góp của đồ án là xây dựng kiến trúc phần mềm thống nhất, kết hợp điều khiển thủ công, điều khiển tự động, giám sát thời gian thực, quản lý dữ liệu nhiệm vụ và nhận diện vật thể, qua đó hỗ trợ vận hành ROV trực quan, an toàn và hiệu quả hơn.


\begin{flushright}
Sinh viên thực hiện\\
\begin{tabular}{@{}c@{}}
\textit{(Ký và ghi rõ họ tên)}
\end{tabular}
\end{flushright}

\end{document}